% TOP_PROBLEM: {"id": 20000936, "title": "A linear worst-case lower bound for Ramsey numbers of rectilinear knots", "category_id": 2, "category": {"id": 2, "name": "combinatorics", "display_name": "Combinatorics", "description": "Counting problems, graph theory, discrete structures.", "slug": "combinatorics", "order_index": 2, "created_at": "2026-07-31T15:26:25.671Z"}, "status": "partially_solved", "problem_number": "AIM-COMBINATORICS-0061", "set_id": 14, "statusQualification": "Uses the source ambient-homeomorphism containment convention, meaning the knot or its mirror; explains the distinct isotopy-only variant."}
% FIRST_POSTED: 2026-10-01
% ENTRY_KIND: statement_only
% UnsolvedMath ID 20000936; source catalogue code AIM-COMBINATORICS-0061
% !TeX program = lualatex
% Definition and sources only; this file is not an LLM solution attempt.
\documentclass[11pt,a4paper]{article}
\usepackage[margin=25mm]{geometry}
\usepackage{fontspec}
\setmainfont{lmroman10-regular.otf}[BoldFont=lmroman10-bold.otf,
 ItalicFont=lmroman10-italic.otf,BoldItalicFont=lmroman10-bolditalic.otf]
\usepackage{amsmath,amssymb,mathtools}
\usepackage{unicode-math}
\setmathfont{latinmodern-math.otf}
\AtBeginDocument{\let\setminus\smallsetminus}
\usepackage{xurl,hyperref}
\hypersetup{colorlinks=true,urlcolor=blue}
\setcounter{secnumdepth}{0}
\setlength{\parindent}{0pt}
\setlength{\parskip}{5pt}
\setlength{\emergencystretch}{3em}
\begin{document}
\section*{A linear worst-case lower bound for Ramsey numbers of rectilinear knots}\label{20000936}
{\small UnsolvedMath ID 20000936 · AIM-COMBINATORICS-0061 · Combinatorics · AIM Workshop Problem Lists}
\subsection{Definitions and mathematical statement}
A tame knot is an embedded circle in $\mathbb R^3$ considered up to ambient isotopy.
Its crossing number $\operatorname{cr}(K)$ is the smallest number of crossings in any regular planar knot diagram.
A rectilinear embedding of the complete graph $K_N$ places $N$ distinct vertices in $\mathbb R^3$ and realizes each edge by its straight segment, with distinct edges meeting only at their common endpoints.
Each graph cycle is consequently a polygonal knot.

Let $R_{\rm knot}(K)$ be the least $N$ such that every such embedding contains a cycle ambient isotopic to $K$ or to its mirror image.
This is the source's ambient-homeomorphism convention for containment, allowing a reflection of three-dimensional space.
Define
\[
 R_{\rm knot}(k)=\max_{\operatorname{cr}(K)\leq k}R_{\rm knot}(K).
\]
The maximum is over the finitely many knot types with at most $k$ crossings, and is well-defined using the knot-containment theorem quoted in the workshop.
Determine the growth of $R_{\rm knot}(k)$ as $k\to\infty$, with upper bounds valid for every geometric placement and lower bounds furnished by placements avoiding a prescribed knot type.
Containment requires a single cycle in the embedded graph, not a drawing obtained after moving its vertices.
The oriented ambient-isotopy variant can also be considered by forbidding the mirror alternative; its convention must be stated separately.
The central issue is quantitative dependence on knot complexity, rather than the qualitative existence of a finite threshold.
\subsection{Short English statement}
How large must a straight-line complete graph in space be to force every knot of a prescribed crossing complexity?
\subsection{Sources}
\begin{itemize}
\item[S1] ULAM.AI and the UnsolvedMath Contributors, supplied problems.json, record 20000936 (AIM-COMBINATORICS-0061), AIM Workshop Problem Lists. Catalogue identity and source transcription; CC BY 4.0.
\url{https://huggingface.co/datasets/ulamai/UnsolvedMath}.
\item[S2] American Institute of Mathematics, Graph Ramsey theory, item 14. Original workshop question underlying AIM-COMBINATORICS-0061.
\url{https://www.overleaf.com/read/mnvcscjjysvg}.
\end{itemize}
\end{document}
